\chapter{Comptes} \label{accounts-chapter}


\vspace{-0.5cm}
\noindent
\cjdcipersonname{COUTANT}{Thomas} \
\cjdcipersonname{LITAMPHA}{Benoît} \
\cjdcipersonname{GEHANT}{Aurélie}
% Alternance des lignes : blanc / gris clair

%===========TABLEAU===========
\rowcolors{2}{white}{gray!15}
\begin{table}[h]
    \centering
    \begin{small}
    \pgfplotstabletypeset[
        col sep=comma,
        string type,
        header=true,
        every head row/.style={
            before row=\hline\rowcolor{gray!50}, after row=\hline
        },
        every even row/.style={
            before row=\rowcolor{gray!30}
        },
        every row/.style={after row=\hline},
        every last row/.style={after row=\hline}, % <-- ligne à ajouter
        columns/Recette/.style={
            column name=Recettes, column type=|p{5cm}
        },
        columns/MontantRecette/.style={
            column name=Montant (€), column type=r||
        },
        columns/Dépense/.style={
            column name=Dépenses, column type=p{5cm}
        },
        columns/MontantDépense/.style={
            column name=Montant (€), column type=r|
        }
    ]{files/accounts.csv}
    \end{small}
    \caption{Table des recettes et dépenses de la compagnie CJDCI (2024–2025)}
\end{table}

%===========GRAPHIQUE===========
\newcommand{\drawBar}[4]{%
  \ifx#3\myempty
    % Si valeur vide : ne rien dessiner
  \else
    \ifdim #3pt = 0pt
      % Si valeur nulle : ne rien dessiner
    \else
      \ifdim #3pt > 0pt
        % Barre positive (recette)
        \draw[fill=green!60] (#1,0) rectangle +(0.8,#3);
        \node[above] at (#1+0.4,#3) {\small #4};
        \node[below=10pt, rotate=90, anchor=east] at (#1+0.4,0) {\scriptsize #2};
      \else
        % Barre négative (dépense)
        \draw[fill=red!60] (#1,0) rectangle +(0.8,#3);
        \node[below] at (#1+0.4,#3) {\small #4};
        \node[above=10pt, rotate=90, anchor=west] at (#1+0.4,0) {\scriptsize #2};
      \fi
    \fi
  \fi
}
    

\newcommand{\drawBarChart}{%
  \begin{tikzpicture}[x=2cm,y=0.02cm]
    \draw[->] (-0.5,0) -- (8.5,0);     % Axe horizontal
    \draw[->] (0,-160) -- (0,260);     % Axe vertical

    \foreach \i in {0,...,7} {
      \pgfmathtruncatemacro{\idx}{\i+1}

      \edef\label{\directlua{
        local data = prepare_data()
        tex.sprint(data[\idx][1] or "")
      }}

      \edef\scaledval{\directlua{
        local data = prepare_data()
        tex.sprint(data[\idx][2] or 0)
      }}

      \edef\realval{\directlua{
        local data = prepare_data()
        tex.sprint(data[\idx][3] or 0)
      }}

      \drawBar{\i}{\label}{\scaledval}{\realval}
    }
  \end{tikzpicture}
}

\begin{luacode*}
function prepare_data()
  -- Fichier CSV contenant les données : 2 colonnes pour les recettes, 2 pour les dépenses
  local filename = "files/accounts.csv"

  -- Tableaux pour stocker les recettes et les dépenses séparément
  local recettes = {}
  local depenses = {}

  -- Lecture du fichier ligne par ligne
  for line in io.lines(filename) do
    -- Extraction des 4 champs (séparés par des virgules) : label et montant pour recette et dépense
    local rlabel, rmontant, dlabel, dmontant = line:match("^([^,]*),([^,]*),([^,]*),([^,]*)")

    -- Conversion des montants texte en nombres (si échec => 0)
    local rval = tonumber(rmontant) or 0
    local dval = tonumber(dmontant) or 0

    -- Si la ligne de recette est valide (non vide, pas un en-tête ou total, et valeur non nulle)
    if rlabel and rlabel ~= "Recette" and rlabel ~= "Total" and rval ~= 0 then
      table.insert(recettes, {rlabel, rval, rval}) -- on stocke : nom, valeur mise à l’échelle (plus tard), valeur réelle
    end

    -- Idem pour les dépenses (on stocke la valeur négative pour affichage en bas du graphe)
    if dlabel and dlabel ~= "Dépense" and dlabel ~= "Total" and dval ~= 0 then
      table.insert(depenses, {dlabel, -dval, -dval}) -- valeurs négatives pour barre descendante
    end
  end

  -- On complète à 4 éléments si besoin (pour éviter les index manquants dans le graphique)
  while #recettes < 4 do table.insert(recettes, {"", 0, 0}) end
  while #depenses < 4 do table.insert(depenses, {"", 0, 0}) end

  -- Fusion des deux tableaux dans un seul pour traitement et affichage
  local all = {}
  for i=1,4 do all[i] = recettes[i] end
  for i=1,4 do all[i+4] = depenses[i] end

  -- Recherche de la valeur absolue maximale pour l’échelle du graphique
  local max_val = 0
  for i,v in ipairs(all) do
    local abs_v = math.abs(v[2]) -- valeur absolue
    if abs_v > max_val then max_val = abs_v end
  end

  -- Calcul du facteur d’échelle pour adapter à la hauteur max du graphe (100 unités visuelles)
  local scale = max_val / 100
  if scale == 0 then scale = 1 end -- éviter division par zéro

  -- Mise à l’échelle de toutes les valeurs à dessiner
  for i,v in ipairs(all) do
    v[2] = v[2] / scale -- on conserve la valeur réelle en v[3]
  end

  -- Retourne le tableau des données prêtes à dessiner
  return all
end
\end{luacode*}



\begin{figure}[H]
  \centering
  \drawBarChart
    \caption{Diagramme des recettes et dépenses de la compagnie CJDCI (2024–2025)}
  \label{fig:recettes-depenses}
\end{figure}


\end{document}
